% Section 5: sparse linear algebra.
\section{Sparse linear algebra}\label{sec:linalg}

Every engine except PDLP spends most of its time solving linear systems: the simplex with the basis
matrix $B$, the LP interior-point method with the normal matrix $M=B\Theta B\T$, and the QP method with
the augmented matrix $K$. \nirnay{} implements the three factorisations it needs and the
fill-reducing ordering they share. All numeric kernels are Numba-compiled loops over plain CSC
arrays.

\subsection{Sparse LU of the simplex basis}\label{sec:lu}
\file{linalg/lu.py} factorises $BQ=\tilde L U$, where $Q$ is a column order, $U$ is upper
triangular in pivot order, and $\tilde L$ is unit lower triangular up to a row permutation (column
$k$ of $\tilde L$ has its unit entry at the pivot row $p_k$).

\paragraph{Column order}
A simplex basis is mostly triangular, and the order exploits that (Suhl and Suhl~\cite{suhl1990}).
A linear-time pass over the basis repeatedly removes \emph{column singletons} of the active
submatrix (logical columns are the first of them), then \emph{row singletons}; each removal can
create new singletons. What is left is the \emph{bump}, ordered by increasing active column count.
Each triangular column is given the singleton row found for it as its pivot, and the factorisation
takes the columns in the order column singletons, row singletons (in the order found), bump. In
the left-looking scheme below, a row-singleton row is touched by no later column, so its column of
$\tilde L$ is never applied again. All fill is therefore confined to the bump. On final Netlib bases
this lowers $(\nnz(\tilde L)+\nnz(U))/\nnz(B)$ from 2.95 to 1.81 on \code{d2q06c}, 2.47 to 1.80 on
\code{25fv47} and 1.24 to 1.06 on \code{degen3}; \code{pilot}, whose bump is large and dense, stays
near 7 and needs a genuine Markowitz factorisation of the bump.

\paragraph{Left-looking elimination}
Column $k$ is computed by the sparse triangular solve of Gilbert and Peierls~\cite{gilbert1988}:
\begin{equation}\label{eq:gp}
  \tilde L_{:,1:k-1}\, x = B_{:,q_k},
\end{equation}
where only columns of $\tilde L$ already computed take part. The rows that become nonzero in
$x$ are exactly the nodes reachable from the nonzeros of $B_{:,q_k}$ in the graph of $\tilde L$. An
iterative depth-first search finds them in topological order, so the numeric solve touches only
those rows. The work is proportional to the floating-point operations, not to~$m$. Entries of $x$
in already-pivoted rows form column $k$ of~$U$; the others are the candidates for the pivot.

\paragraph{Threshold pivoting}
A triangular column takes its assigned singleton row whenever that entry passes the threshold test.
Otherwise, among unpivoted rows $r$ with $|x_r|\ge\tau\max_{r'}|x_{r'}|$ ($\tau=0.1$), the pivot is
the row with the fewest nonzeros inside the bump, with ties broken by magnitude. This is Markowitz's idea
\cite{markowitz1957} applied row-wise: it gives up a little stability for much less fill.

\paragraph{Basis repair}
If no candidate exceeds $10^{-11}$, column $q_k$ depends on the columns before it. Rather than
reject the basis, the factorisation leaves step $k$ unpivoted and reports a pair (basis position,
unpivoted row). The simplex then swaps that basic variable for the row's logical, which is always
independent of the rest, and refactorises. Near-singular bases arise from rounding and from
degenerate pivots on models such as the \code{pilot} family, and industrial simplex codes keep
going the same way.

\paragraph{Updates}
After a pivot that replaces basis position $r$ by a column with $\alpha_q=B^{-1}a_q$, the new inverse
is $B_{\text{new}}^{-1}=E^{-1}B^{-1}$ with the elementary eta matrix
\begin{equation}\label{eq:eta}
  E = I + (\alpha_q - e_r)e_r\T,\qquad
  E^{-1}v = v - \frac{v_r}{\alpha_{q,r}}(\alpha_q - e_r) ,
\end{equation}
the product form of the inverse~\cite{dantzig1954}. The eta file stores $\alpha_q$ sparsely,
dropping entries below $10^{-14}$. FTRAN ($B x = b$) applies the LU solve and then the etas in order;
BTRAN ($B\T y=d$) applies the etas in reverse and then the transposed LU solve. For that
solve the factorisation also keeps row-wise copies of $\tilde L$ and $U$, so both triangular solves
run in the axpy form that skips zero entries; the right-hand side $e_r$ of the simplex's BTRAN
is very sparse, and the column-wise dot-product form could not exploit that. The basis is
refactorised after 100 updates, or earlier when the eta file holds more than
$8(\nnz(\tilde L)+\nnz(U)+m)$ entries, to bound both the work per solve and the growth of rounding
error.

\begin{algorithm}[t]
\caption{Left-looking sparse LU with threshold pivoting and repair (\file{linalg/lu.py})}\label{alg:lu}
\Input{basis $B$ (CSC), column order $q$, row counts $\rho_r$, threshold $\tau$, tiny $\epsilon$}
\Output{$\tilde L$, $U$, pivot rows $p$, list of unpivoted steps}
\For{$k=1,\dots,m$}{
  $\mathcal{X}\leftarrow$ reach of $\operatorname{struct}(B_{:,q_k})$ in the graph of $\tilde L$, in topological order (DFS)\;
  scatter $x\leftarrow B_{:,q_k}$; \ForEach{$r\in\mathcal X$ with $r$ pivoted at step $s$}{$x_{\mathcal{L}_s}\leftarrow x_{\mathcal{L}_s}-\tilde L_{\mathcal{L}_s,s}\,x_r$}
  $U_{:,k}\leftarrow\{x_r : r \text{ pivoted}\}$;\quad $a_{\max}\leftarrow\max\{|x_r|: r\text{ unpivoted}\}$\;
  \eIf{$a_{\max}>\epsilon$}{
    $p_k\leftarrow\arg\min\{\rho_r : r \text{ unpivoted},\ |x_r|\ge\tau a_{\max}\}$ \Comment*[r]{shortest row}
    $U_{kk}\leftarrow x_{p_k}$;\quad $\tilde L_{:,k}\leftarrow x_{\text{unpivoted}}/x_{p_k}$\;
  }{
    record step $k$ as singular \Comment*[r]{caller swaps in a logical}
  }
  clear $x$ on $\mathcal X$\;
}
\end{algorithm}

\subsection{Sparse Cholesky for the normal equations}\label{sec:chol}
The LP interior-point method solves $M\,\Delta y=r$ with $M=B\Theta B\T+\delta I$ symmetric
positive definite and the same sparsity pattern at every iteration. \file{linalg/cholesky.py}
splits the work in the classic way~\cite{davis2006}:
\begin{enumerate}
  \item \textbf{Analyse} once. Compute a fill-reducing permutation $P$ (\cref{sec:ordering}), the
        upper triangle of $PMP\T$ with a map back to $M$'s values, the elimination tree
        ($\operatorname{parent}(i)=\min\{k>i: l_{ki}\ne0\}$, computed with path compression), and
        the column counts of $L$ from the row subtrees. $L$'s pattern is then allocated exactly.
  \item \textbf{Factor} every iteration, by the up-looking algorithm. Row $k$ of $L$ solves
        \begin{equation}\label{eq:uplook}
          L_{1:k-1,1:k-1}\; l_{k} = m_{1:k-1,k},\qquad
          L_{kk} = \sqrt{m_{kk} - l_k\T l_k},
        \end{equation}
        and the nonzeros of $l_k$ are the nodes on the paths from the nonzeros of $m_{1:k-1,k}$
        up the elimination tree to $k$ (the row subtree).
  \item \textbf{Solve} $LL\T\hat x = P b$ by forward and back substitution.
\end{enumerate}

\paragraph{Pivot guard}
As the interior-point iterates converge, $M$ becomes nearly singular: rows may be dependent,
and $\Theta$ spans many orders of magnitude. If the pivot $m_{kk}-l_k\T l_k\le 10^{-30}$, it is
replaced by $10^{128}$. The corresponding component of the solution then becomes (almost) zero
instead of overflowing, which is the standard remedy~\cite[§11.1]{wright1997}. The number of
replaced pivots is reported with each factorisation.

\subsection{Minimum-degree ordering and fill-in}\label{sec:ordering}
Eliminating a node of the graph of a symmetric matrix joins all its neighbours into a clique, and
the new edges are exactly the fill-in of~$L$. \file{linalg/ordering.py} implements the minimum-degree
heuristic~\cite{tinney1967,george1989}: it repeatedly eliminates a node of smallest current degree,
on an explicit elimination graph with degree buckets. Two refinements matter on optimisation
matrices:
\begin{itemize}
  \item \textbf{mass elimination:} neighbours whose closed neighbourhood equals the pivot's would each
        be chosen next anyway, so they are eliminated together. This removes most of the cost on
        block-structured matrices;
  \item \textbf{dense-row deferral:} nodes with degree above $\max(16,10\sqrt n)$ are taken out of the
        graph and ordered last. Normal matrices of LPs often have a few almost-dense rows, and
        ordering them last both bounds the work of the ordering and keeps the dense part of $L$ in
        the final block.
\end{itemize}
\Cref{fig:fill} shows why the order matters on the smallest example, an arrowhead matrix.

\begin{figure}[tbp]
\centering
\newcommand{\cell}[3]{\fill[#3] ({(#2-1)*0.42},{-(#1-1)*0.42}) rectangle ++(0.36,-0.36);}
\begin{tikzpicture}
  % ---- left: hub first
  \begin{scope}
    \draw[rule] (-0.06,0.06) rectangle (2.52,-2.52);
    \foreach \i in {1,...,6}{\cell{\i}{\i}{navy}}
    \foreach \j in {2,...,6}{\cell{1}{\j}{navy}\cell{\j}{1}{navy}}
    \foreach \i in {2,...,6}{\foreach \j in {2,...,6}{\ifnum\i=\j\else\cell{\i}{\j}{saffron}\fi}}
    \node[lbl,text=ink,align=center] at (1.23,-2.95) {hub eliminated first\\[-1pt]20 fill entries};
  \end{scope}
  % graph in the middle
  \begin{scope}[shift={(4.4,-1.23)}]
    \node[circle,draw=navy,fill=navy,text=white,inner sep=1.2pt,font=\sffamily\scriptsize] (h) at (0,0) {1};
    \foreach \k/\a in {2/90,3/162,4/234,5/306,6/18}{
      \node[circle,draw=navy,fill=white,inner sep=1.2pt,font=\sffamily\scriptsize] (n\k) at (\a:1.0) {\k};
      \draw[navy] (h) -- (n\k);}
    \draw[saffron,densely dashed] (n2)--(n3) (n3)--(n4) (n4)--(n5) (n5)--(n6) (n6)--(n2) (n2)--(n4) (n2)--(n5) (n3)--(n5) (n3)--(n6) (n4)--(n6);
    \node[lbl,text=ink,align=center] at (0,-1.72) {graph of $M$; dashed: clique\\[-1pt]created by eliminating node 1};
  \end{scope}
  % ---- right: hub last (minimum degree)
  \begin{scope}[shift={(6.9,0)}]
    \draw[rule] (-0.06,0.06) rectangle (2.52,-2.52);
    \foreach \i in {1,...,6}{\cell{\i}{\i}{navy}}
    \foreach \j in {1,...,5}{\cell{6}{\j}{navy}\cell{\j}{6}{navy}}
    \node[lbl,text=ink,align=center] at (1.23,-2.95) {minimum degree: hub last\\[-1pt]no fill};
  \end{scope}
  \node[lbl,anchor=west] at (9.7,-0.2) {\tikz\fill[navy] (0,0) rectangle (0.25,0.25); nonzero of $M$};
  \node[lbl,anchor=west] at (9.7,-0.7) {\tikz\fill[saffron] (0,0) rectangle (0.25,0.25); fill in $L+L\T$};
\end{tikzpicture}
\caption[Fill-in and ordering]{Fill-in and ordering on an arrowhead matrix. Eliminating the hub
first (left) makes its five neighbours a clique (middle), and $L$ becomes dense. Minimum degree
eliminates the leaves first and the hub last (right), and $L$ has no fill at all. The same effect,
on a large scale, decides whether the normal matrix of an LP can be factorised at all.}
\label{fig:fill}
\end{figure}

\subsection{Assembling the normal matrix}\label{sec:normal}
$M = B\Theta B\T + \diag(\delta)$ changes its values at every interior-point iteration but not its
pattern. \file{linalg/normal.py} therefore computes once the full pattern of $BB\T$ (with sorted
row indices) and a \emph{slot map}. For each row $k$ of $B$, each column $j$ in that row, and each
entry $i$ of column $j$, the map gives the position of $M_{ik}$ in the value array. Each
iteration then assembles $M$ in one pass,
\begin{equation}\label{eq:assemble}
  M_{\text{slot}(q)} \mathrel{+}= \theta_j\, b_{kj}\, b_{ij},\qquad q=1,2,\dots,\ \sum_j \nnz(b_j)^2 ,
\end{equation}
with no searching and no allocation, followed by the diagonal term $\delta_k$.

\subsection{Quasi-definite \texorpdfstring{$LDL\T$}{LDLT} for the QP augmented system}\label{sec:ldl}
The QP method solves systems with the regularised augmented matrix
\begin{equation}\label{eq:qd}
  K=\begin{bmatrix} -H & B\T\\ B & R\end{bmatrix},\qquad H = Q+D+\rho I\succ0,\quad R=\delta I\succ0,
\end{equation}
which is symmetric \emph{quasi-definite}. Vanderbei~\cite{vanderbei1995} showed that such a matrix
has a factorisation $PKP\T=LDL\T$ with $D$ diagonal, $n$ negative and $m$ positive pivots, for
\emph{every} symmetric permutation $P$. The ordering can therefore be chosen for sparsity alone,
exactly as for Cholesky, and no pivoting for stability is needed during the numeric
factorisation. \file{linalg/ldl.py} reuses the Cholesky analysis (minimum degree, elimination tree,
column counts) and the up-looking $LDL\T$ kernel of Davis~\cite{davis2005}. The numeric step computes,
for row $k$,
\begin{equation}\label{eq:ldl}
  L_{kj} = \frac{y_j}{D_{jj}},\qquad D_{kk} = k_{kk} - \sum_{j<k} L_{kj}\,y_j ,
\end{equation}
where $y$ solves the sparse triangular system for row $k$.

\paragraph{Pivot floor}
The signs of the pivots are known in advance: $\sigma_k=-1$ for a primal row of $K$ and $+1$ for a
dual row. A pivot with the wrong sign, or smaller than $\delta_{\min}$ in magnitude, is replaced by
$\sigma_k\delta_{\min}$ (dynamic regularisation, with $\delta_{\min}=10^{-11}$ in the QP method). The
factorisation then always exists and has the correct inertia. The perturbation is removed by
iterative refinement against the unregularised matrix (\cref{sec:qp}). \Cref{fig:kkt} shows the
block structure.

\begin{figure}[tbp]
\centering
\begin{tikzpicture}[x=1cm,y=1cm]
  % K blocks
  \fill[saffron!22] (0,0) rectangle (3.2,-3.2);
  \fill[navy!14] (3.2,0) rectangle (5.2,-3.2);
  \fill[navy!14] (0,-3.2) rectangle (3.2,-5.2);
  \fill[igreen!18] (3.2,-3.2) rectangle (5.2,-5.2);
  \draw[navy,line width=0.7pt] (0,0) rectangle (5.2,-5.2);
  \draw[navy,line width=0.4pt] (3.2,0) -- (3.2,-5.2) (0,-3.2) -- (5.2,-3.2);
  % sparsity motif
  \foreach \p in {(0.3,-0.3),(0.7,-0.7),(1.1,-1.1),(1.5,-1.5),(1.9,-1.9),(2.3,-2.3),(2.7,-2.7),(3.05,-3.05),
                  (0.7,-1.5),(1.5,-0.7),(1.1,-2.3),(2.3,-1.1)}{\fill[saffron!85!black] \p circle (0.05);}
  \foreach \p in {(3.5,-3.5),(3.9,-3.9),(4.3,-4.3),(4.7,-4.7),(5.0,-5.0)}{\fill[igreen] \p circle (0.05);}
  \foreach \p in {(3.5,-0.5),(3.5,-1.4),(3.9,-0.9),(4.3,-2.1),(4.7,-2.6),(5.0,-1.8),(4.3,-0.4),
                  (0.5,-3.5),(1.4,-3.5),(0.9,-3.9),(2.1,-4.3),(2.6,-4.7),(1.8,-5.0),(0.4,-4.3)}{\fill[navy] \p circle (0.05);}
  \node[font=\small,text=ink,fill=saffron!22,inner sep=2pt] at (1.6,-1.6) {$-(Q+D+\rho I)$};
  \node[font=\small,text=ink,fill=navy!14,inner sep=2pt] at (4.2,-1.6) {$B\T$};
  \node[font=\small,text=ink,fill=navy!14,inner sep=2pt] at (1.6,-4.2) {$B$};
  \node[font=\small,text=ink,fill=igreen!18,inner sep=2pt] at (4.2,-4.2) {$\delta I$};
  % dimension braces
  \draw[decorate,decoration={brace,amplitude=4pt},muted] (0,0.12) -- node[above=3pt,lbl]{$n_B$ primal} (3.2,0.12);
  \draw[decorate,decoration={brace,amplitude=4pt},muted] (3.2,0.12) -- node[above=3pt,lbl]{$m_B$ dual} (5.2,0.12);
  % solution / rhs
  \draw[navy] (6.0,0) rectangle (6.5,-5.2); \draw[navy] (6.0,-3.2) -- (6.5,-3.2);
  \node[font=\small] at (6.25,-1.6) {$\Delta x$}; \node[font=\small] at (6.25,-4.2) {$\Delta y$};
  \node[font=\small] at (6.9,-2.6) {$=$};
  \draw[navy] (7.3,0) rectangle (7.8,-5.2); \draw[navy] (7.3,-3.2) -- (7.8,-3.2);
  \node[font=\small] at (7.55,-1.6) {$\hat r$}; \node[font=\small] at (7.55,-4.2) {$r_b$};
  % D signs
  \begin{scope}[shift={(9.0,0)}]
    \node[lbl,anchor=south] at (0.25,0.05) {pivots of $D$};
    \foreach \k in {0,...,7}{\fill[saffron] (0,-\k*0.4) rectangle (0.5,-\k*0.4-0.36); \node[font=\sffamily\scriptsize,text=white] at (0.25,-\k*0.4-0.18) {$-$};}
    \foreach \k in {0,...,4}{\fill[igreen] (0,-3.2-\k*0.4) rectangle (0.5,-3.2-\k*0.4-0.36); \node[font=\sffamily\scriptsize,text=white] at (0.25,-3.2-\k*0.4-0.18) {$+$};}
    \node[lbl,align=left,anchor=west] at (0.7,-1.6) {$n_B$ negative};
    \node[lbl,align=left,anchor=west] at (0.7,-4.2) {$m_B$ positive};
  \end{scope}
\end{tikzpicture}
\caption[Quasi-definite KKT block structure]{Block structure of the regularised augmented
(KKT) system solved at every QP interior-point iteration. The $(1,1)$ block is negative definite,
because it contains the barrier term $D=Z X^{-1}+S W^{-1}$ and the regularisation $\rho I$, and the
$(2,2)$ block is positive definite. Any symmetric permutation therefore factorises as $LDL\T$ with
the pivot signs shown on the right, and the ordering can be chosen for sparsity alone.}
\label{fig:kkt}
\end{figure}
